\documentclass[11pt]{article}
\usepackage[margin=1in]{geometry}
\usepackage{amsmath,amssymb,booktabs,array,longtable}
\usepackage{float}
\usepackage[T1]{fontenc}
\usepackage{lmodern}
\usepackage{microtype}
\usepackage{hyperref}
\hypersetup{colorlinks=true,linkcolor=blue,urlcolor=blue}
\newcommand{\TopK}{\operatorname{TopK}}
\newcommand{\KL}{\operatorname{KL}}
\newcommand{\sig}{\operatorname{sigmoid}}
\newcommand{\loss}{\mathcal{L}}
\title{Functional Task Networks: Approximately Bayesian State Space Routing with Input-Conditioned Priors}
\author{Kevin L. McKee}
\date{}
\begin{document}
\maketitle

\noindent A functional task network (FTN) is a masked model of selectable
functional units; its low-rank adaptation (LoRA) variant is FTN-LoRA.
This document specifies the feed-forward implementations in
\texttt{projects/bayesian\_adv}. Unless stated otherwise, ``default'' means the
current \texttt{Bayes}, \texttt{BayesLoRA}, or composable
\texttt{InputConditioned} configuration in \texttt{config.py}, rather than a
constructor fallback. Let $H$ be the number of selectable units, $K$ the
number active in one hard mask, $S=\{(x_i,y_i)\}_{i=1}^n$ a support batch, and
$\loss_S(f)=n^{-1}\sum_i\ell(f(x_i),y_i)$. The implementation uses cross-entropy
for classification and mean squared error for regression.

\section{Background: FTN and FTN-LoRA}
\paragraph{FTN.}
An FTN has $H$ independently parameterized small multilayer perceptrons
(MLPs) $u_j(x;\theta_j)$. Their scalar outputs use rectified linear unit
(ReLU) hidden layers and a final $\tanh$; a shared
readout $W\in\mathbb{R}^{C\times H}$ combines them. With a binary mask
$m\in\{0,1\}^H$, $\sum_jm_j=K$, its prediction is
\begin{equation}
 z_j(x)=u_j(x;\theta_j),\qquad
 f_{\theta,W}(x;m)=W\bigl(m\odot z(x)\bigr).
 \label{eq:ftn}
\end{equation}
Dropout, when enabled, acts on $z$ before masking. In the ordinary FTN,
mask reconfiguration searches temporary real-valued logits against $\loss_S$
with model weights detached, then installs a hard top-$K$ mask. The selected
weights are subsequently trained on data under that mask. Thus the mask
selects which independent functional units contribute to the readout; it is
not a separate learned network for every task.

\paragraph{FTN-LoRA.}
FTN-LoRA instead has one frozen random MLP. Each linear layer $\ell$ has a
bank of $H$ trainable rank-$r$ atoms, and the \emph{same} mask selects atoms
at every layer. Using row-vector notation, the layer computes
\begin{equation}
 h_{\ell+1}=h_\ell W^0_\ell+
 \frac{\alpha}{K}\sum_{j=1}^{H}m_j
 \left[(\operatorname{Drop}_\ell(h_\ell)A_{\ell j})B_{\ell j}
       +b_{\ell j}\right],
 \label{eq:lora}
\end{equation}
followed by ReLU except at the output layer. Here $W^0_\ell$ is frozen,
$A_{\ell j}\in\mathbb{R}^{d_\ell\times r}$,
$B_{\ell j}\in\mathbb{R}^{r\times d_{\ell+1}}$, and the optional
$b_{\ell j}$ is a trainable per-atom bias. There is no frozen base bias.
The bank is evaluated densely and multiplied by $m$, so all atoms have a
mask-gradient path during straight-through search even when their current
forward mask value is zero. Only atoms and their biases receive ordinary
weight updates. Equations~\eqref{eq:ftn} and \eqref{eq:lora} share a $K$-of-$H$
router, but the former gates outputs of independent MLPs while the latter
gates edits to a shared computation.

\section{Bayes FTN: persistent mask state}
Bayes FTN retains the FTN predictor in Eq.~\eqref{eq:ftn}. Its additional
state is a logit vector $L_t\in\mathbb{R}^{H}$, initialized to zero and
carried across training batches. $L_t$ is a buffer updated by mask inference,
not a parameter of the weight optimizer. For an ordinary update, the detached
prior logits and the starting optimization variable are
\begin{equation}
 P_t=\rho L_{t-1},\qquad \lambda^{(0)}=P_t.
 \label{eq:prior}
\end{equation}
$\rho=1$ preserves the preceding state exactly; $\rho=0$ gives zero logits,
whose Bernoulli surrogate probabilities are uniform at $1/2$. We call this
a \emph{persistent-state prior}: it is the preceding approximate mask state,
possibly diffused toward zero. No other prior construction is needed here.
With zero logits, Gumbel top-$K$ is uniform over exact-$K$ subsets, so each
unit's actual inclusion probability is $K/H$, not the surrogate's $1/2$.
For the configured masks, $K/H=1/16$ in Bayes FTN and $1/4$ in BayesLoRA.
This illustrates why the Bernoulli probabilities below are an optimization
proxy rather than calibrated marginals of the hard-mask distribution.

At each of $J$ adaptive moment estimation (Adam) steps, independent
$g_j=-\log(-\log U_j)$ with
$U_j\sim\mathrm{Uniform}(0,1)$ perturb the ranking. The forward mask is
exactly $K$-sparse:
\begin{align}
 a_s&=\exp\!\left[-\frac{s}{\max\{(J-1)/3,10\}}\right],
 &m^{(s)}&=\TopK_K\!\left(\frac{\lambda^{(s)}}{T a_s}+g\right),\nonumber\\
 q_j&=\sig(\lambda_j/T),
 &p_j&=\sig(P_{t,j}/T).
 \label{eq:gumbel}
\end{align}
When stochastic sampling is disabled, the forward mask is simply
$\TopK_K(\lambda^{(s)})$. The search uses the straight-through quantity
$\widetilde m=m^{(s)}+q-\operatorname{stopgrad}(q)$: its numerical forward
value is hard, while its derivative follows the sigmoid relaxation. Backbone
weights are detached during this search. A Kullback--Leibler (KL) divergence
penalizes movement away from the prior. The objective is
\begin{equation}
 \min_{\lambda}\quad
 \loss_S\bigl(f_{\theta,W}(\cdot;\widetilde m)\bigr)
 +\beta\,\frac{1}{H}\sum_{j=1}^{H}
 \left[q_j\log\frac{q_j}{p_j}
 +(1-q_j)\log\frac{1-q_j}{1-p_j}\right].
 \label{eq:objective}
\end{equation}
This is an \emph{approximate} Bayesian update. The mean factorized-Bernoulli
KL is an analytic, coordinate-separable $O(H)$ penalty with a deterministic
gradient; it uses the same sigmoid probabilities as the straight-through
path. These properties make it inexpensive during the few Adam steps of an
ordinary update. It is a stability penalty toward prior logits, not the KL
of the exact-$K$ distribution used in forward passes. In particular,
Gumbel top-$K$ couples selections through sampling without replacement,
whereas the Bernoulli surrogate treats coordinates independently. It also
penalizes a common shift $\lambda\mapsto\lambda+c\mathbf{1}$ even though
that shift leaves the top-$K$ law unchanged.

More distribution-matched alternatives are possible. The \emph{ordered}
Gumbel top-$K$ draw $i_1,\ldots,i_K$ has a Plackett--Luce probability.
Writing $a_j=\lambda_j/T_s$ for the sampling temperature at a given step,
\begin{equation}
 \log Q_a(i_1,\ldots,i_K)=\sum_{r=1}^{K}
 \left[a_{i_r}-\log\!\sum_{j\notin\{i_1,\ldots,i_{r-1}\}}e^{a_j}\right].
 \label{eq:pl}
\end{equation}
Its log probability is easy to evaluate, so sampled differences between
the current and prior log probabilities could estimate a KL over \emph{orders}.
The expectation and its gradient would need Monte Carlo estimation, however,
and the order contains information discarded by the FTN's binary mask.
Consequently, the ordered KL is an upper bound on the corresponding
unordered-mask KL, not that KL itself. Obtaining an unordered mask's exact
probability requires accounting for its possible selection orders.
Another choice is a conditional-Bernoulli exact-$K$ distribution,
$Q(S)\propto\exp(\sum_{j\in S}\lambda_j/T)$ for $|S|=K$.
Its normalization can be computed by dynamic programming in $O(HK)$ rather
than enumerating all $\binom{H}{K}$ masks, but it is a different sampling law
from Gumbel top-$K$ and is costlier than the present $O(H)$ penalty.

At the end, the optimized logits become $L_t$ and a hard mask becomes the
active $m_t$.
With \texttt{TAKE\_BEST=True}, Bayes FTN instead retains the sampled mask
and pre-step logits that attained the lowest support loss during search;
the selection criterion is task loss, not the KL-augmented objective.
$L_t$ carries information to the next mask update; $m_t$ determines the
present forward computation and which unit parameters receive task gradients.
The active hard mask, rather than a fresh sample on every prediction, is used
for subsequent weight updates until reconfiguration. The reactive detector
below uses an exponential moving average (EMA) of recent losses; configuration
keys abbreviate learning rate as \texttt{LR} and standard deviation as
\texttt{SD}.

\begin{table}[H]
\centering\small
\caption{Bayes FTN defaults from \texttt{Bayes} (and inherited \texttt{Default}).
Sample sizes are rows of the current support batch.}
\begin{tabular}{@{}p{.40\linewidth}p{.17\linewidth}p{.34\linewidth}@{}}
\toprule
Configuration key & Default & Role \\
\midrule
\texttt{SQUARE\_DIM}, $H$ & $64$, $4096$ & Unit-bank size ($H=64^2$) \\
\texttt{K} & $256$ & Active units per mask \\
\texttt{LAYERS}, \texttt{INNER\_DIM} & $8$, $8$ & Depth and hidden width of each unit MLP \\
\texttt{BAYES\_RECONFIG\_STEPS} & $2$ & Ordinary posterior steps $J$ \\
\texttt{BAYES\_RECONFIG\_LR} & $0.3$ & Adam rate for logits \\
\texttt{RECONFIG\_SAMPLE\_SIZE} & $256$ & Ordinary support limit \\
\texttt{BAYES\_KL\_WEIGHT} & $1.0$ & $\beta$ in Eq.~\eqref{eq:objective} \\
\texttt{BAYES\_TEMPERATURE} & $3.0$ & $T$ for sigmoid and sampling \\
\texttt{BAYES\_PRIOR\_DECAY} & $1.0$ & $\rho$ in Eq.~\eqref{eq:prior} \\
\texttt{BAYES\_STOCHASTIC\_MASK} & \texttt{True} & Gumbel perturbations on \\
\texttt{BAYES\_RECONFIG\_TAKE\_BEST} & \texttt{True} & Retain lowest support-loss iterate \\
\texttt{RECONFIG\_REACTIVE}, \texttt{RECONFIG\_REACTIVE\_SD}, \texttt{RECONFIG\_REACTIVE\_EMA} & \texttt{True}, $5.0$, $0.9$ & Loss-triggered schedule; threshold and EMA factor \\
\texttt{EVAL\_RECONFIG\_STEPS}, \texttt{EVAL\_RECONFIG\_RESTARTS} & $30$, $1$ & Independent cold evaluation search \\
\bottomrule
\end{tabular}
\end{table}

Evaluation-time \texttt{cold\_reconfigure} is separate from training-state
updates: it starts temporary logits at zero, optimizes on labeled support
without a KL term, and selects the best hard mask by support loss. It does
not modify $L_t$. Normal evaluation uses 30 steps, Adam rate $0.3$,
temperature $1$, one restart, and the full supplied support unless a sample
limit is configured.

\section{BayesLoRA}
BayesLoRA applies the same persistent-logit update,
Eqs.~\eqref{eq:prior}--\eqref{eq:objective}, to the FTN-LoRA predictor of
Eq.~\eqref{eq:lora}. Its $H$ mask coordinates now select low-rank atoms,
not independent output units. During mask inference, all trainable atom
weights and biases are detached; during the following weight update, the
installed hard mask selects their contributions. This distinction matters:
overlapping masks expose shared atom parameters to updates from multiple
solutions, whereas nonoverlapping atoms are structurally isolated. The
frozen base matrices are shared but never updated.

\begin{table}[H]
\centering\small
\caption{BayesLoRA defaults that differ from Bayes FTN, plus its LoRA-specific
backbone values. Unlisted posterior and detector settings are inherited from
Table 1.}
\begin{tabular}{@{}p{.40\linewidth}p{.17\linewidth}p{.34\linewidth}@{}}
\toprule
Configuration key & Default & Role \\
\midrule
\texttt{K} & $1024$ & Active atoms from $H=4096$ \\
\texttt{LORA\_LAYERS}, \texttt{LORA\_WIDTH} & $8$, $256$ & Frozen MLP hidden depth and width \\
\texttt{LORA\_RANK} & $1$ & Rank of each atom \\
\texttt{LORA\_ALPHA} & $256$ & Scale $\alpha/K$ of summed edits \\
\texttt{LORA\_WEIGHT\_SCALE} & $1.0$ & Frozen base initialization scale \\
\texttt{LORA\_DROPOUT}, \texttt{LORA\_BIAS} & $0$, \texttt{True} & Atom-input dropout; bias bank \\
\texttt{LORA\_B0\_SCALE} & $0.01$ & Initial $B$ and bias standard deviation \\
\texttt{BAYES\_RECONFIG\_TAKE\_BEST} & \texttt{False} & Install final logits and draw final mask \\
\texttt{TRAIN\_DURATION} & $1500$ & Training batches per task block (Bayes: $500$) \\
\texttt{REPLAY\_BATCH\_SIZE} & $16$ & Stored rows per replay batch (Bayes: $256$) \\
\bottomrule
\end{tabular}
\end{table}

\section{Changepoint detection and cold reset/refit}
The reactive training loop tests each incoming batch \emph{before} changing
its mask or weights. Let $d_t=\loss_{S_t}(f(\cdot;m_{t-1}))$, let $\mu_t$ and
$v_t$ be the running post-update loss mean and variance, and let $c_t$ count
batches since the last detected change. Apart from the first batch, it flags
a change when
\begin{equation}
 c_t\geq \left\lceil\frac{1}{1-\gamma}\right\rceil
 \quad\text{and}\quad
 d_t>\mu_{t-1}+s\sqrt{\max(v_{t-1},0)}.
 \label{eq:detector}
\end{equation}
Here $\gamma=0.9$, $s=5$, and the warmup is ten batches. After fitting the
mask, the loss $a_t$ under the new mask restarts $\mu=a_t,v=0$ on a change;
otherwise the code updates
$\delta=a_t-\mu$, $\mu\leftarrow\gamma\mu+(1-\gamma)a_t$,
$v\leftarrow\gamma v+(1-\gamma)\delta^2$. The detector never reads a task
identifier or declared task boundary. Between detections, BayesLoRA still
makes its two-step persistent update on each batch.

On detection, \texttt{reconfigure\_changepoint} uses the same mask-inference
objective with a separate parameter set. Its defaults are
\begin{equation}
 J_{\rm cp}=30,\quad \eta_{\rm cp}=0.3,\quad
 \rho_{\rm cp}=0,\quad\beta_{\rm cp}=0,\quad T_{\rm cp}=1,
 \label{eq:cp}
\end{equation}
with stochastic Gumbel masks and the entire detected batch as support
(the changepoint sample-size option is \texttt{None}). Hence
$P_t=0$ and $\lambda^{(0)}=0$: this is the long reconfiguration from
uniform prior logits, or \emph{cold start}, rather than a continuation of
the old mask state. The resulting logits become the new persistent state.
BayesLoRA keeps the final logits and draws its installed mask from them.
This training-time reset/refit is distinct from the stateless evaluation
cold search described above.

\clearpage
\section{Input-conditioned priors}
Abstractly, input-conditioned routing maps a small input set $X_s$ to
mask scores $h(X_s)\in\mathbb{R}^H$. The useful choice is what those scores
\emph{do}: they may determine a mask directly, or provide initial logits and
a prior for the existing support-set optimization,
\begin{equation}
 \begin{aligned}
 L_0&=h(X_s),\\
 L^*&\approx\operatorname{Optimize}_{L\,\leftarrow\,L_0}
 \left[\loss_S\bigl(f(\cdot;\widetilde m(L))\bigr)
       +\beta_{\rm ic}\KL(q_L\Vert q_{L_0})\right],\\
 m&=\TopK_K(L^*).
 \end{aligned}
 \label{eq:ic-abstract}
\end{equation}
Here $S$ adds labeled support when refinement is requested; the hard mask
and its straight-through search follow Section 2. Thus input conditioning
need not replace mask inference. It can supply a better starting point while
the task loss corrects an imperfect match. Direct routing is the choice to
use $m=\TopK_K(L_0)$ without that refinement.

A learned neural map $h$ would have its own trainable parameters. Sequential
updates for new input domains could then catastrophically forget old
data-to-mask associations, causing routing failure even if the masked FTN
still retains those solutions. The implemented input-conditioned (IC)
mechanism avoids additional trainable router weights: it identifies a domain
using analytically fitted, fixed input projectors and retrieves saved mask
logits. This protects the \emph{routing map} from gradient-driven forgetting;
it does not by itself prevent retrieved logits from becoming stale when the
FTN's shared trainable parameters change. The repair option below addresses
that separate problem.

Concretely, the router adds an archive of pairs $(P_a,L_a)$ to Bayes FTN or
BayesLoRA, where $P_a$ is an input-domain projector and $L_a$ is a saved
persistent mask-logit vector. For the current segment, it accumulates the
\emph{uncentered} scatter
$C=\sum_i x_i^\top x_i$ over unlabeled training inputs. If $U_r$ contains
the top $r$ eigenvectors of $C$, it builds the rank-$r$ projector
$P=U_rU_r^\top$. A detected changepoint archives the outgoing projector
and logits and begins collecting a new segment. Archived entries are
append-only; a current-segment candidate may also participate in evaluation.

Given the first $n_s$ unlabeled rows $X_s$ of a support batch, route matching
uses normalized reconstruction error
\begin{equation}
 e(P;X_s)=\frac{\|X_s-X_sP\|_F^2}{\max(\|X_s\|_F^2,\varepsilon)}.
 \label{eq:reconstruction}
\end{equation}
The router takes the \emph{latest} entry whose error is within tolerance
$\delta$ of the minimum, accepting it only if its error is below threshold
$\tau$. The mapping $h(X_s)=L_a$ therefore retrieves an
\emph{input-conditioned prior} over mask choices. Matching uses inputs alone;
labels are involved only if a subsequent refinement or fallback search is
requested. Input distributions
that are indistinguishable by Eq.~\eqref{eq:reconstruction} cannot be
distinguished by this router merely because their labels differ.

At evaluation, a successful match has two selectable behaviors.
\texttt{direct} applies $\TopK_K(L_a)$ immediately, requiring neither labels
nor optimization. The configured default \texttt{refine} starts temporary
logits at $L_a$ and performs $J_{\rm ic}$ deterministic-mask posterior steps
on labeled support, with $\rho=1$ and KL weight $\beta_{\rm ic}$ toward
$L_a$. It installs the resulting hard mask but restores the pre-evaluation
persistent logits. Thus retrieval and support-based optimization are
complementary in the default path. A match failure takes a third path:
ordinary zero-logit cold search on labeled support. With the current
defaults, refinement uses ten steps and at most four labeled rows; fallback
uses the ordinary 30-step search with the same four-row limit.

Training-time use is a separate switch. With the default
\texttt{IC\_TRAIN\_PRIOR=False}, detected changepoints retain the long
zero-logit reset/refit of Eq.~\eqref{eq:cp}. If enabled, a matched
changepoint instead initializes from $L_a$ and uses $J_{\rm ic}$ steps with
$\rho=1$ and $\beta_{\rm ic}$; a miss still cold-starts. Another independent
switch, \texttt{IC\_NOVELTY\_TRIGGER}, can add a changepoint when the live
segment projector reconstructs current inputs with error at least $\tau$;
it is off by default. The ordinary loss detector remains active either way.

Saved logits may become stale when shared model weights change. The optional,
default-enabled replay repair operates after each training block: it groups stored
labeled replay examples by archived input projector, cold-searches a mask on
fit rows, and accepts it only if a separate validation set improves by more
than the specified margin over the old mask. Acceptance \emph{appends} a new
entry with the old projector and logits $+a$ on selected atoms and $-a$
elsewhere; it does not edit earlier entries or add trainable router weights.
Disabling \texttt{IC\_REPLAY\_REPAIR} leaves archived masks fixed, which can
allow their predictions to degrade as shared FTN weights evolve. Repair uses
labeled training replay even though domain matching itself uses unlabeled
inputs; it therefore requires sufficient stored rows for its separate fit
and validation sets.

\begin{longtable}{@{}p{.40\linewidth}p{.16\linewidth}p{.35\linewidth}@{}}
\caption{Defaults in the composable \texttt{InputConditioned} configuration.
The inherited Bayes/BayesLoRA values remain as above. In configuration keys,
\texttt{AE} abbreviates autoencoder; here its fitted model is a linear projector.}\\
\toprule
Configuration key & Default & Role \\
\midrule
\endfirsthead
\toprule Configuration key & Default & Role \\ \midrule
\endhead
\texttt{IC\_SUPPORT\_SIZE} & $4$ & Unlabeled rows used in matching \\
\texttt{IC\_AE\_RANK} & $2$ & Projector rank $r$ \\
\texttt{IC\_RECON\_THRESHOLD} & $0.10$ & Route acceptance threshold $\tau$ \\
\texttt{IC\_MATCH\_TOLERANCE} & $10^{-3}$ & Latest-near-best tolerance $\delta$ \\
\texttt{IC\_SAME\_DOMAIN\_DISTANCE} & $0.05$ & Projector-distance threshold for repair deduplication \\
\texttt{IC\_NOVELTY\_TRIGGER} & \texttt{False} & Add live-domain novelty to detector \\
\texttt{IC\_TRAIN\_PRIOR} & \texttt{False} & Use retrieved logits at training changes \\
\texttt{IC\_EVAL\_STRATEGY} & \texttt{refine} & Direct mask or labeled refinement \\
\texttt{IC\_PRIOR\_STEPS} & $10$ & $J_{\rm ic}$ for prior refinement \\
\texttt{IC\_PRIOR\_KL} & $1.0$ & $\beta_{\rm ic}$ toward retrieved logits \\
\texttt{EVAL\_RECONFIG\_SAMPLE\_SIZE} & $4$ & Labeled rows for refinement or fallback \\
\texttt{IC\_REPLAY\_REPAIR} & \texttt{True} & Enable retrospective mask repair \\
\texttt{IC\_REPAIR\_EVERY\_BLOCKS} & $1$ & Repair interval \\
\texttt{IC\_REPAIR\_FIT\_ROWS}, \texttt{IC\_REPAIR\_VALIDATION\_ROWS} & $800$, $800$ & Repair fit and held-out rows \\
\texttt{IC\_REPAIR\_STEPS}, \texttt{IC\_REPAIR\_LR} & $100$, $0.3$ & Cold-search steps and Adam rate \\
\texttt{IC\_REPAIR\_TEMPERATURE}, \texttt{IC\_REPAIR\_RESTARTS} & $1.0$, $3$ & Cold-search sampling and restarts \\
\texttt{IC\_REPAIR\_TAKE\_BEST} & \texttt{True} & Retain best search candidate \\
\texttt{IC\_REPAIR\_MIN\_IMPROVEMENT} & $0.0$ & Required validation-loss gain \\
\texttt{IC\_REPAIR\_LOGIT\_MAGNITUDE} & $4.0$ & Accepted entry's logit magnitude $a$ \\
\bottomrule
\end{longtable}

\appendix
\section{FTN as a state-space router}
An architecture-independent FTN can be written as a state-space router over
a bank of functional units. Let $s_t$ be routing state,
$\theta_t=(\theta_{t1},\ldots,\theta_{tH})$ the current trainable unit
parameters, and $B_t$ the available support. A schematic cycle is
\begin{align}
 s_t&=\mathcal{R}(s_{t-1},\theta_t,B_t),
 &m_t&=\TopK_K\bigl(\mathcal{S}(s_t)\bigr),\nonumber\\
 \hat y_t&=f_{\theta_t}(x_t;m_t)
 =\mathcal{C}\!\left(x_t;
      m_{t1}\mathcal{U}_1(\theta_{t1}),\ldots,
      m_{tH}\mathcal{U}_H(\theta_{tH}),\theta_0\right),\nonumber\\
 \theta_{t+1}&=\theta_t-\eta\nabla_{\theta_t}
 \loss_{B_t}\bigl(f_{\theta_t}(\cdot;m_t)\bigr).
 \label{eq:router}
\end{align}
$\theta_0$ denotes shared fixed structure. The last line is an idealized
gradient step; the implementation uses Adam for weight training. These
updates are coupled: routing chooses the trainable units exposed on this
batch, and their changed parameters alter the fit of future routes.
$\mathcal{U}_j$ is a functional unit and $\mathcal{C}$ combines the selected
units. The route state need
not be probabilistic: it could hold deterministic search scores, a mask,
or a memory index. In the feed-forward FTN, $\mathcal{U}_j$ is a small MLP's
scalar output and $\mathcal{C}$ is a linear readout. In FTN-LoRA,
$\mathcal{U}_j$ is the collection of atom edits
$\{A_{\ell j}B_{\ell j},b_{\ell j}\}_{\ell}$ and $\mathcal{C}$ is the
composition of the edited layers with nonlinearities.

The central dynamical property is conditional parameter isolation. For a
hard route, an atom or independent MLP absent from $m_t$ contributes nothing
to the forward calculation and receives no ordinary task-loss gradient:
\begin{align}
 m_{tj}=0&\quad\Longrightarrow\quad
 \nabla_{\theta_j}\loss_{B_t}=0,\nonumber\\
 \operatorname{supp}(m_t)\cap\operatorname{supp}(m_{t'})
 &\quad\text{identifies parameters exposed to both updates.}
 \label{eq:isolation}
\end{align}
This statement concerns the ordinary task-loss gradient; optimizer momentum
can still move an inactive parameter after it has received earlier gradients.
It assumes each $\theta_j$ belongs only to unit $j$; FTN's
shared readout is best viewed by its per-unit columns, while the LoRA base
is fixed. The mask intersection identifies where interference or transfer
\emph{can} occur, not their magnitude or sign: those depend on the data,
loss, and parameter gradients. The router can change overlap over time by
reusing, dropping, or recruiting units.
Thus isolation is \emph{dynamical}, rather than a once-for-all partition of
parameters among tasks.

\subsection*{A Bayesian representation}
The same router admits a conditional latent-variable description. Let
$\mathcal{M}_K=\{m\in\{0,1\}^H:\sum_jm_j=K\}$ and let $Q_L$ denote the
distribution on $\mathcal{M}_K$ induced by Gumbel top-$K$ with logits $L$
at a fixed sampling temperature.
The previous routing state defines a predictive prior, and a single mask
explains the labeled support batch:
\begin{align}
 \pi_t(m)&=Q_{\rho L_{t-1}}(m),\qquad m\in\mathcal{M}_K,\nonumber\\
 p_{\theta_t}(Y_t\mid X_t,m)
   &=\prod_{(x_i,y_i)\in B_t}p_{\theta_t}(y_i\mid x_i,m),\nonumber\\
 p(m\mid B_t,L_{t-1},\theta_t)
   &=\frac{p_{\theta_t}(Y_t\mid X_t,m)\,\pi_t(m)}
     {\sum_{m'\in\mathcal{M}_K}
       p_{\theta_t}(Y_t\mid X_t,m')\,\pi_t(m')}.
 \label{eq:bayes-router}
\end{align}
For classification, $p_{\theta_t}(y\mid x,m)$ can be the categorical
distribution obtained by applying softmax to the masked network's outputs;
for squared-error regression, a Gaussian with masked-network mean and fixed
variance gives a corresponding likelihood. Only the route is treated as
latent here; the functional-unit weights $\theta_t$ are point estimates.

An ideal variational router would approximate the posterior in
Eq.~\eqref{eq:bayes-router} by a tractable $Q_L$:
\begin{equation}
 L_t^*\in\arg\min_L\left\{
 \mathbb{E}_{m\sim Q_L}\!\left[-\log p_{\theta_t}(Y_t\mid X_t,m)\right]
 +\KL\!\left(Q_L\,\Vert\,\pi_t\right)\right\},
 \qquad m_t\sim Q_{L_t^*}.
 \label{eq:variational-router}
\end{equation}
This equation specifies an exact-$K$ Bayesian target, not the implemented
optimizer. The implementation carries $L_t$ forward but replaces the
more involved exact-mask KL with the factorized-Bernoulli
surrogate of Eq.~\eqref{eq:objective}, estimates task fit with hard-mask
samples and straight-through gradients, and takes only a few steps. Its
mean loss and separately weighted mean KL should therefore be read as a
practical approximate update, not a calibrated posterior computation.
At a changepoint, setting $\rho=0$ resets $\pi_t$ to a uniform exact-$K$
prior; the implemented refit also sets the KL weight to zero, making that
search an unregularized fit. An input-conditioned prior can instead set prior logits from a
retrieved $h(X_s)$ before optional support-based refinement. In all cases,
the probabilistic route changes which parameters are exposed to learning;
the isolation and overlap property in Eq.~\eqref{eq:isolation} does not
depend on using a Bayesian router.

\end{document}
